%! Author  = Omar Iskandarani
%! Date    = June 2026
%! ORCID   = 0009-0006-1686-3961
%! Target  = Physical Review A / European Physical Journal Plus

% ============================================================
%  A Compton-Core Parametrization of Vortex-Filament Scales:
%  Mechanical Scales from a Single Geometric Postulate and
%  the Spectroscopic Rydberg Anchor
% ============================================================

\documentclass[aps,pra,preprint,longbibliography,nofootinbib]{revtex4-2}

\usepackage{amsmath,amssymb,amsthm,mathtools}
\usepackage{booktabs}
\usepackage{enumitem}
\usepackage{siunitx}
\usepackage[colorlinks=true,
            linkcolor=blue!60!black,
            citecolor=blue!60!black,
            urlcolor=blue!60!black]{hyperref}
\usepackage{cleveref}
\usepackage{microtype}

% ── Theorem environments ─────────────────────────────────────
\newtheorem{postulate}{Postulate}
\newtheorem{lemma}{Lemma}
\newtheorem{proposition}{Proposition}
\newtheorem{corollary}{Corollary}
\newtheorem{remark}{Remark}

% ── Notation ─────────────────────────────────────────────────
% Circulatory arrow subscript (reads as "circulatory flow")
\newcommand{\rotarrow}{%
  \mathchoice{\mkern-2mu\scriptstyle\boldsymbol{\circlearrowleft}}%
             {\mkern-2mu\scriptstyle\boldsymbol{\circlearrowleft}}%
             {\mkern-2mu\scriptscriptstyle\boldsymbol{\circlearrowleft}}%
             {\mkern-2mu\scriptscriptstyle\boldsymbol{\circlearrowleft}}}

\newcommand{\vc}{\mathbf{v}_{\rotarrow}}        % circulatory velocity (vector)
\newcommand{\vnorm}{\lVert\mathbf{v}_{\!\boldsymbol{\circlearrowleft}}\rVert} % magnitude

\newcommand{\rhof}{\rho_{\!f}}                   % background fluid density
\newcommand{\rhoE}{\rho_{\!E}}                   % kinetic energy density
\newcommand{\rhom}{\rho_{\!m}}                   % mass-equivalent density
\newcommand{\rhocore}{\rho_{\mathrm{core}}}      % core density
\newcommand{\rc}{r_c}                            % core radius
\newcommand{\Stv}{S_t}                           % relativistic time-dilation factor
\newcommand{\omC}{\omega_C}                      % Compton angular frequency
\newcommand{\lc}{\lambda_c}                      % Compton wavelength
\newcommand{\Rinf}{R_\infty}                     % Rydberg constant
\newcommand{\Tc}{\mathcal{T}_c}                  % vortex line tension at core
\newcommand{\Fgr}{F_{\mathrm{gr}}}              % gravitational tension scale c^4/4G
\newcommand{\Gz}{\Gamma_0}                       % circulation quantum
% ─────────────────────────────────────────────────────────────

\begin{document}

\title{A Compton-Core Parametrization of Vortex-Filament Scales:\\
       Mechanical Scales from a Single Geometric Postulate\\
       and the Spectroscopic Rydberg Anchor}

\author{Omar Iskandarani}
\affiliation{Independent Researcher, Groningen, The Netherlands}
\email{info@omariskandarani.com}
\orcid{0009-0006-1686-3961}
\date{June 2026}

% ─────────────────────────────────────────────────────────────
\begin{abstract}
We present a self-consistent parametrization of the mechanical scales
of a localized vortex-core model of the electron from a single
structural postulate combined with the spectroscopic Rydberg constant
$\Rinf$.
The postulate identifies the angular rotation rate of the vortex core
with the electron Compton frequency $\omC = m_e c^2/\hbar$.
Together with $\Rinf$, which is algebraically equivalent to the
fine-structure constant $\alpha$ given $\{m_e, c, \hbar\}$,
this yields the characteristic circulation speed
$\vnorm = (\pi\hbar c\Rinf/m_e)^{1/2}$
and anchors the complete set of vortex-core scales.
The core radius $\rc = \vnorm/\omC = r_e/2$ (half the classical
electron radius), the vortex line tension
$\Tc = m_e c^2/(2\rc)$ (Compton energy per core-diameter length),
the core mass density $\rhocore = m_e c^2/(2\pi\vnorm^2\rc^3)$
from a Laplace-pressure balance, and the circulation quantum
$\Gz = 2\pi\vnorm^2/\omC$ follow as algebraic consequences,
each with a direct classical-hydrodynamic interpretation.
A status table makes explicit which quantities are postulated,
which are calibrated, and which are derived.
Two parameter-free falsifiable predictions are stated.
\end{abstract}

\maketitle

% =============================================================
\section{Introduction}
\label{sec:intro}
% =============================================================

Localized vortex-filament models of fundamental particles trace their
lineage to the proposals of Helmholtz and
Kelvin~\cite{Helmholtz1867,Kelvin1867} and have received renewed
attention in the contexts of topological solitons~\cite{FaddeevNiemi1997},
analogue gravity~\cite{Barcelo2005}, and superfluid vacuum
models~\cite{Volovik2003}.
A characteristic circulation speed $\vnorm$ and a core radius $\rc$
appear as fundamental scales in all such frameworks, but their
numerical values are typically supplied by fitting to known particle
masses rather than derived from independent constraints.

The present paper asks: given a single geometric postulate about the
core's rotation rate, what set of mechanical scales is implied?
We show that the postulate, together with the well-measured Rydberg
constant $\Rinf$, fixes all commonly used vortex-core parameters in a
non-circular algebraic chain.
The parameter $\Rinf$ is informationally equivalent to the
fine-structure constant $\alpha$ via $\Rinf = \alpha^2 m_e c/(2h)$;
choosing $\Rinf$ as the anchor is a matter of experimental access,
since $\Rinf$ carries the lowest relative uncertainty of any
combination of $\{\alpha, m_e, c, h\}$.

The contribution of this paper is not the derivation of $\alpha$ from
independent physics — that would require going outside standard
quantum electrodynamics — but the explicit demonstration that the
Compton-core postulate provides a geometrically and
hydrodynamically motivated single-parameter closure for the
vortex-core model.
All derived quantities have direct classical-hydrodynamic
interpretations in terms of solid-body rotation, circulation,
line tension, and Laplace-pressure balance.

Throughout we work in SI units, with $\hbar = h/(2\pi)$, $m_e$,
$c$, $e$ carrying their CODATA 2018 meanings~\cite{MohrTaylor2018}.
We denote the circulatory velocity field by $\vc$ and its magnitude
by $\vnorm = |\vc|$.

% =============================================================
\section{The Compton-Core Postulate}
\label{sec:postulate}
% =============================================================

\begin{postulate}[Compton-core hypothesis]
\label{post:compton}
The vortex core undergoes solid-body rotation at the electron
Compton angular frequency,
\begin{equation}
  \Omega_{\mathrm{core}} = \omC \equiv \frac{m_e c^2}{\hbar}.
  \label{eq:postulate}
\end{equation}
The magnitude of the circulatory velocity at the core boundary
of radius $\rc$ is therefore
\begin{equation}
  \vnorm = \omC\,\rc.
  \label{eq:vc_rc_link}
\end{equation}
\end{postulate}

\begin{remark}[Physical motivation]
\label{rem:motivation}
For a solid-body rotating core of radius $\rc$ and angular
velocity $\Omega$, the tangential speed at the boundary is
$\vnorm = \Omega\,\rc$.
The Compton frequency $\omC = m_e c^2/\hbar$ is the natural
frequency associated with the electron rest energy: it is the
frequency at which one quantum $\hbar\omC = m_e c^2$ is stored.
Postulate~\ref{post:compton} is the simplest resonance condition
consistent with the electron rest energy as the core-energy scale,
and it provides a geometric link between the rotation rate and the
core boundary.
It is a hypothesis about the core dynamics, not a consequence of
standard quantum electrodynamics, and it is testable: any
independent measurement of $\rc$ inconsistent with
$\rc = \vnorm/\omC$ (where $\vnorm$ is determined by
Sec.~\ref{sec:derived}) falsifies the postulate.
\end{remark}

% =============================================================
\section{The Rydberg Anchor}
\label{sec:rydberg}
% =============================================================

\subsection{Equivalence of $\Rinf$ and $\alpha$ as inputs}

The Rydberg constant satisfies the exact relation
\begin{equation}
  \Rinf = \frac{\alpha^2 m_e c}{2h},
  \label{eq:Rydberg_standard}
\end{equation}
so that $\Rinf$ and $\alpha$ carry identical information given
$\{m_e, c, h\}$.
We use $\Rinf$ rather than $\alpha$ as the anchor for two reasons:
(a) $\Rinf$ has the smallest relative uncertainty,
$\delta\Rinf/\Rinf = \SI{1.9e-12}{}$, of any combination of
$\{\alpha, m_e, c, h\}$~\cite{MohrTaylor2018};
(b) it is determined from hydrogen spectroscopy independently of
assumptions about the electron's internal structure.

\subsection{Rydberg identity in core variables}

\begin{lemma}[Rydberg identity]
\label{lem:rydberg}
Given Postulate~\ref{post:compton}, the Rydberg constant satisfies
\begin{equation}
  \Rinf = \frac{\vnorm^2\,\omC}{\pi c^3}
         = \frac{\vnorm^3}{\pi\,\rc\,c^3}.
  \label{eq:rydberg_core}
\end{equation}
This is an algebraic identity: it holds as a consequence of the
standard definition~\eqref{eq:Rydberg_standard} and the relation
$\alpha = 2\vnorm/c$ (established in Sec.~\ref{sec:alpha}).
Its utility here is as a \emph{closure relation}: given only
$\{m_e, c, \hbar, \Rinf\}$ and Postulate~\ref{post:compton},
it uniquely determines $\vnorm$ without prior knowledge of $\alpha$.
\end{lemma}

\begin{proof}
Substitute $\rc = \vnorm/\omC$ into the right-hand side:
\[
  \frac{\vnorm^3}{\pi(\vnorm/\omC)c^3}
  = \frac{\vnorm^2\omC}{\pi c^3}
  = \frac{\vnorm^2 m_e}{\pi\hbar c}. 
\]
Set $\vnorm = \alpha c/2$:
\[
  = \frac{\alpha^2 m_e c}{4\pi\hbar}
  = \frac{\alpha^2 m_e c}{2h}
  = \Rinf. \qed
\]
\end{proof}

\begin{remark}
\label{rem:algebraic_status}
The identity~\eqref{eq:rydberg_core} does not derive $\alpha$ from
independent physics: it is an algebraic reformulation of
$\Rinf = \alpha^2 m_e c/(2h)$ in terms of $\vnorm$ and $\rc$.
The non-trivial content lies in Postulate~\ref{post:compton}, which
supplies the constraint $\rc = \vnorm/\omC$.
Together with~\eqref{eq:rydberg_core}, this constraint determines
$\vnorm$ uniquely from $\Rinf$.
\end{remark}

% =============================================================
\section{Derivation of the Characteristic Speed}
\label{sec:derivation}
% =============================================================

Substituting $\rc = \vnorm/\omC$ into Lemma~\ref{lem:rydberg}:
\begin{equation}
  \Rinf = \frac{\vnorm^2\,\omC}{\pi c^3}
  \quad\Longrightarrow\quad
  \vnorm^2 = \frac{\pi\,\hbar\,c\,\Rinf}{m_e}.
  \label{eq:vc_solve}
\end{equation}

\begin{proposition}[Characteristic circulation speed]
\label{prop:vc}
Given Postulate~\ref{post:compton} and the measured $\Rinf$,
\begin{equation}
  \boxed{
    \vnorm = \sqrt{\frac{\pi\,\hbar\,c\,\Rinf}{m_e}}
  }
  \label{eq:vc}
\end{equation}
is the unique positive closure value for the tangential
circulation speed.
The inputs are $\{m_e, c, \hbar, \Rinf\}$; since
$\Rinf \equiv \alpha^2 m_e c/(2h)$, this is equivalent to using
$\alpha$ as input.
\end{proposition}

\subsubsection*{Dimensional check}
\[
  \Bigl[\sqrt{\hbar\,c\,\Rinf/m_e}\Bigr]
  = \sqrt{
      \frac{\mathrm{J{\cdot}s\cdot m{\cdot}s^{-1}\cdot m^{-1}}}
           {\mathrm{kg}}
    }
  = \sqrt{\mathrm{m^2\,s^{-2}}} = \mathrm{m\,s^{-1}}. \checkmark
\]

\subsubsection*{Numerical value}
$\vnorm = \SI{1.093\,845\,631\,5e6}{\meter\per\second}$
\quad(from CODATA 2018~\cite{MohrTaylor2018}).

% =============================================================
\section{Derived Mechanical Scales}
\label{sec:derived}
% =============================================================

The following scales follow algebraically from
Proposition~\ref{prop:vc} and Postulate~\ref{post:compton},
each with a direct classical-hydrodynamic interpretation.

% ------- 4.1 Core radius ------------------------------------
\subsection{Core radius and its relation to the classical electron radius}
\label{sec:rc}

\begin{equation}
  \rc = \frac{\vnorm}{\omC} = \frac{\vnorm\,\hbar}{m_e c^2}
      = \SI{1.408\,970\,16e-15}{\meter}.
  \label{eq:rc}
\end{equation}

\begin{remark}[Relation to the classical electron radius]
\label{rem:re}
The classical electron radius is
\begin{equation}
  r_e = \frac{e^2}{4\pi\varepsilon_0 m_e c^2}
      = \frac{\alpha\hbar}{m_e c}
      = \SI{2.817\,940\,3e-15}{\meter}.
  \label{eq:re}
\end{equation}
From \eqref{eq:rc} and $\alpha = 2\vnorm/c$ (established below),
$\rc = \alpha\hbar/(2m_e c) = r_e/2$.
We emphasize that $\rc$ is an \emph{effective model parameter},
not a claim about the physical extent of the electron.
The classical radius $r_e$ itself plays an analogous role in QED
— appearing in Thomson cross-sections, Compton scattering, and
vacuum polarization — without implying a physical "size."
Experimental upper bounds on the electron's spatial extent
(from $g{-}2$ and deep-inelastic measurements,
$r < \SI{e-18}{\meter}$~\cite{PDG2022}) constrain the
\emph{point-like} quantum-mechanical extent of the charge
distribution, not the scale at which an effective vortex-core
model is constructed.
For a toroidal vortex in which the tube centre lies at distance
$R = \rc$ from the symmetry axis and the tube has minor radius
$a = \rc$ (horn-torus geometry), the outer boundary lies at
$R + a = 2\rc = r_e$, providing a geometric interpretation of
the factor-of-two relation.
\end{remark}

% ------- 4.2 Fine-structure constant -------------------------
\subsection{Fine-structure constant as algebraic equivalence}
\label{sec:alpha}

\begin{proposition}[Algebraic equivalence of $\alpha$ and $\vnorm$]
\label{prop:alpha}
The ratio $\tilde\alpha \equiv 2\vnorm/c$
satisfies $\tilde\alpha = \alpha$ (the CODATA fine-structure
constant) to the precision of the input data.
This is not a derivation of $\alpha$ from independent physics
but the algebraic consequence of $\Rinf \equiv \alpha^2 m_e c/(2h)$
combined with \eqref{eq:vc}.
\end{proposition}

\begin{proof}
$\tilde\alpha = (2/c)(\pi\hbar c\Rinf/m_e)^{1/2}$.
Substituting $\Rinf = \alpha^2 m_e c/(4\pi\hbar)$:
$\tilde\alpha = (2/c)(\alpha^2 c^2/4)^{1/2} = \alpha$. \qed
\end{proof}

Numerically, $\tilde\alpha = \num{7.297\,352\,567\,1e-3}$
versus CODATA $\alpha = \num{7.297\,352\,569\,3e-3}$;
the residual $3\times10^{-10}$ is the propagated uncertainty
from $\Rinf$ (relative uncertainty \SI{1.9e-12}{}) combined with
the rounding of $\{m_e, c, \hbar\}$.
It is a precision-consistency check, not a prediction.

% ------- 4.3 Circulation quantum ----------------------------
\subsection{Circulation quantum}
\label{sec:Gamma0}

The line integral of velocity around the core boundary defines the
circulation,
\begin{equation}
  \Gz \equiv \oint \vc\cdot d\mathbf{l}
           = 2\pi\,\rc\,\vnorm
           = \frac{2\pi\vnorm^2}{\omC}
           = \SI{9.683\,619e-9}{\meter\squared\per\second}.
  \label{eq:Gamma0}
\end{equation}
The ratio $\Gz/\kappa = \alpha^2/4 \approx \num{1.33e-5}$,
where $\kappa = h/m_e$ is the Onsager-Feynman quantum~\cite{Onsager1949,Feynman1955},
showing that the present circulation is smaller by $\alpha^2/4$.

% ------- 4.4 Vortex line tension ---------------------------
\subsection{Vortex line tension}
\label{sec:Tc}

The \emph{vortex line tension} is the energy stored in the vortex
core per unit length of filament~\cite{Saffman1992}.
We define it from the Compton-energy condition:
\begin{equation}
  \Tc \equiv \frac{m_e c^2}{2\rc} = \frac{\hbar\omC}{2\rc}.
  \label{eq:Tc_def}
\end{equation}
The physical content: the vortex stores one quantum of Compton
energy $m_e c^2 = \hbar\omC$ per segment of filament of length
$2\rc$ (the core diameter).
This is a consequence of Postulate~\ref{post:compton} together
with the requirement that the core energy scale is $m_e c^2$.

\begin{corollary}[Line tension in fundamental constants]
\label{cor:Tc}
\begin{equation}
  \Tc = \frac{m_e c^2}{2\rc}
      = \frac{m_e^2 c^3}{\tilde\alpha\,\hbar}
      = \SI{29.053\,51}{\newton}.
  \label{eq:Tc}
\end{equation}
\end{corollary}

\begin{remark}[Disambiguation from the gravitational tension scale]
\label{rem:Fgr}
The quantity $\Fgr \equiv c^4/(4G)$ defines the maximum force
transmissible without forming a causal horizon in general
relativity~\cite{Gibbons2002,Schiller2006}.
It and $\Tc$ share the dimension of force (energy per unit length),
but are conceptually distinct:
$\Tc$ is a mechanical property of the vortex core, while $\Fgr$ is
a causal upper bound from spacetime physics.
Numerically $\Tc \approx \SI{29}{\newton}$ versus
$\Fgr \approx \SI{3e43}{\newton}$.
Their ratio
\begin{equation}
  \frac{\Tc}{\Fgr} = \frac{4\alpha_g}{\tilde\alpha}
                   \approx \num{9.60e-43},
  \label{eq:tension_ratio}
\end{equation}
where $\alpha_g = Gm_e^2/(\hbar c) \approx \num{1.752e-45}$,
involves only fundamental constants and has no free parameters.
\end{remark}

% ------- 4.5 Core mass density --------------------------------
\subsection{Core mass density from Laplace-pressure balance}
\label{sec:rhocore}

For a cylindrical vortex tube of radius $\rc$ and line tension
$\Tc$, the surface tension of the core wall is
$\gamma = \Tc/(2\pi\rc)$ [N\,m$^{-1}$], and the inward Laplace
pressure is~\cite{Saffman1992,Donnelly1991}
\begin{equation}
  P_L = \frac{\gamma}{\rc} = \frac{\Tc}{2\pi\rc^2}.
  \label{eq:Laplace}
\end{equation}
Setting $P_L$ equal to the dynamic pressure $\tfrac{1}{2}\rhocore\vnorm^2$
of the circulating core fluid:
\begin{equation}
  \frac{1}{2}\rhocore\vnorm^2 = \frac{\Tc}{2\pi\rc^2}.
  \label{eq:pressure_balance}
\end{equation}

\begin{corollary}[Core density]
\label{cor:rho}
Solving~\eqref{eq:pressure_balance} and substituting
\eqref{eq:Tc_def}:
\begin{equation}
  \boxed{
    \rhocore = \frac{\Tc}{\pi\rc^2\vnorm^2}
             = \frac{m_e c^2}{2\pi\,\vnorm^2\,\rc^3}.
  }
  \label{eq:rhocore}
\end{equation}
Numerically, $\rhocore = \SI{3.893\,4e18}{\kilogram\per\cubic\meter}$.
\end{corollary}

\subsubsection*{Dimensional check}
$[m_e c^2/(\vnorm^2\rc^3)]
= \mathrm{kg\,m^2\,s^{-2}\,/\,(m^2\,s^{-2}\cdot m^3)}
= \mathrm{kg\,m^{-3}}$. \checkmark

\begin{remark}[Non-circularity of the Laplace-pressure route]
\label{rem:noncircular}
The derivation of $\rhocore$ via~\eqref{eq:pressure_balance} is
non-circular because $\Tc$ is defined by~\eqref{eq:Tc_def}
independently of $\rhocore$.
Specifically, $\Tc = m_ec^2/(2\rc)$ uses only $\{m_e, c, \rc\}$;
it does not involve $\rhocore$.

For completeness, the Saffman formula for the line tension of a
vortex filament with circulation $\Gz$ and core density
$\rhocore$ is~\cite{Saffman1992}
\begin{equation}
  \Tc^{\mathrm{(Saffman)}} = \frac{\rhocore\,\Gz^2}{4\pi}.
  \label{eq:Tc_saffman}
\end{equation}
Substituting $\Gz = 2\pi\rc\vnorm$ and $\rhocore$
from~\eqref{eq:rhocore}, one verifies
$\Tc^{\mathrm{(Saffman)}} = m_ec^2/(2\rc) = \Tc$ identically.
Equation~\eqref{eq:Tc_saffman} therefore serves as an internal
consistency check rather than an independent definition.
Had $\Tc$ been defined via~\eqref{eq:Tc_saffman} first,
the Laplace balance would have been vacuous (giving $0 = 0$);
the non-trivial content rests on the primary
definition~\eqref{eq:Tc_def}.
\end{remark}

% ------- 4.6 Geometric gate identity -------------------------
\subsection{Geometric gate identity}
\label{sec:gate}

\begin{lemma}[Geometric gate]
\label{lem:gate}
\begin{equation}
  \frac{\lc}{\pi\rc} = \frac{4}{\tilde\alpha} \approx 548.14,
  \label{eq:gate}
\end{equation}
where $\lc = h/(m_ec)$ is the electron Compton wavelength.
\end{lemma}

\begin{proof}
$\lc = 2\pi\hbar/(m_ec)$; $\rc = \tilde\alpha\hbar/(2m_ec)$:
$\lc/(\pi\rc) = 4/\tilde\alpha$. \qed
\end{proof}

This dimensionless ratio is pure algebra; it is a geometric property
of the length scales $\lc$ and $\rc$ and does not encode additional
physical content.

% =============================================================
\section{Derivation Chain and Hydrodynamic Interpretation}
\label{sec:chain}
% =============================================================

The complete derivation chain is:
\begin{equation}
  \underbrace{\omC,\;\Rinf\!\equiv\!\tilde\alpha,\;m_e,\;c,\;\hbar}_{\text{inputs}}
  \;\xrightarrow{\;\text{Post.\;\ref{post:compton}}\;}
  \vnorm
  \xrightarrow{}
  \rc
  \xrightarrow{}
  \Gz
  \xrightarrow{}
  \Tc
  \xrightarrow{}
  \rhocore.
  \label{eq:chain}
\end{equation}
The classical-hydrodynamic interpretation of each scale is:

\medskip
\begin{tabular}{@{}ll@{}}
$\vnorm = \omC\rc$ & tangential speed at solid-body core boundary \\
$\rc = \vnorm/\omC$ & core radius from solid-body rotation \\
$\Gz = 2\pi\rc\vnorm$ & circulation (line integral of velocity) \\
$\Tc = m_e c^2/(2\rc)$ & line tension (Compton energy per core-diameter) \\
$\rhocore$ & core density from Laplace-pressure/dynamic-pressure balance \\
\end{tabular}

% =============================================================
\section{Status Table}
\label{sec:status}
% =============================================================

\begin{table}[h]
\centering
\caption{Epistemic status of all quantities.
``Post.'' = consequence of Postulate~\ref{post:compton};
``Calib.'' = single empirical input ($\Rinf\equiv\tilde\alpha$);
``Algebraic'' = algebraic consequence of prior rows;
``Open'' = not determined within the present framework.}
\label{tab:status}
\begin{tabular}{lllp{4cm}}
\toprule
Quantity & Symbol & Status & Note \\
\midrule
Compton frequency  & $\omC = m_e c^2/\hbar$     & Standard def.  & \cite{MohrTaylor2018} \\
Core rotation rate & $\Omega_{\rm core}=\omC$   & Postulate       & Post.~\ref{post:compton} \\
Rydberg/FSC input  & $\Rinf\equiv\tilde\alpha$  & Calib.\ (input) & $\Rinf=\alpha^2 m_e c/(2h)$ \\
Circ.\ speed       & $\vnorm$                   & Post.+Calib.    & Prop.~\ref{prop:vc} \\
Core radius        & $\rc = r_e/2$              & Post.+Calib.    & \eqref{eq:rc}; effective param. \\
Circulation        & $\Gz$                      & Algebraic       & \eqref{eq:Gamma0} \\
$\tilde\alpha=2\vnorm/c$ & $=\alpha$          & Algebraic equiv.& Prop.~\ref{prop:alpha} \\
Line tension       & $\Tc$                      & Post.+Calib.    & \eqref{eq:Tc_def} \\
Core density       & $\rhocore$                 & Algebraic       & Cor.~\ref{cor:rho}; Laplace \\
Geometric gate     & $\lc/(\pi\rc)=4/\tilde\alpha$ & Algebraic  & Lemma~\ref{lem:gate} \\
Background density & $\rhof$                    & Open            & additional constraint needed \\
Topological state  & knot type $K$              & Open            & not determined \\
\bottomrule
\end{tabular}
\end{table}

% =============================================================
\section{Falsifiable Predictions}
\label{sec:falsifiers}
% =============================================================

\medskip
\noindent\textbf{(i) Core-radius test.}
Postulate~\ref{post:compton} predicts $\rc = \SI{1.408\,970\,16e-15}{\meter}$
as the effective model scale.
Any anomalous short-range contribution to the Coulomb potential at
the femtometre scale detectable in Lamb-shift
measurements~\cite{Bethe1947,Karshenboim2005} that is consistent
with the predicted value would support the model;
absence of any such anomaly at this scale falsifies
Postulate~\ref{post:compton} as a description of the electron.
We note that existing QED calculations account for the Lamb shift
without invoking a core structure at this scale; the prediction
is therefore a bound on new physics rather than a reproduction of
known results.

\medskip
\noindent\textbf{(ii) Line-tension to gravitational-tension ratio.}
Equation~\eqref{eq:tension_ratio} predicts
$\Tc/\Fgr = 4\alpha_g/\tilde\alpha \approx \num{9.60e-43}$
with no fitted parameters.
This dimensionless ratio can in principle be tested in
analogue-gravity systems~\cite{Barcelo2005} where both a mechanical
maximum tension and a causal-horizon analogue are accessible,
though a concrete experimental protocol requires additional modelling
of the analogue system.

% =============================================================
\section{Numerical Summary}
\label{sec:numerical}
% =============================================================

\begin{table}[h]
\centering
\caption{Numerical values of all derived quantities from CODATA 2018
inputs~\cite{MohrTaylor2018}.
The fourth column gives the propagated fractional uncertainty from $\Rinf$
(\num{1.9e-12}) where applicable; dashes indicate quantities with no
CODATA counterpart.}
\label{tab:numbers}
\begin{tabular}{llll}
\toprule
Quantity & Symbol & Value & Prop.\ unc.\ from $\Rinf$ \\
\midrule
Circ.\ speed   & $\vnorm$ & \SI{1.093\,845\,63e6}{\m\per\s}      & $\sim\num{1e-12}$ \\
Core radius    & $\rc$    & \SI{1.408\,970\,16e-15}{\meter}       & $\sim\num{1e-12}$ \\
$\tilde\alpha$ &          & \num{7.297\,352\,567e-3}              & $\sim\num{1e-12}$ \\
Circulation    & $\Gz$    & \SI{9.683\,619e-9}{\m\squared\per\s}  & — \\
Line tension   & $\Tc$    & \SI{29.053\,51}{\newton}               & — \\
Core density   & $\rhocore$& \SI{3.893\,4e18}{\kg\per\m\cubed}    & — \\
Geom.\ gate    & $4/\tilde\alpha$ & \num{548.14}               & $\sim\num{1e-12}$ \\
\bottomrule
\end{tabular}
\end{table}

% =============================================================
\section{Discussion}
\label{sec:discussion}
% =============================================================

The central result of this paper is the explicit parametrization of
all vortex-core mechanical scales from a single geometric postulate
(core rotation at $\omC$) and one spectroscopic input ($\Rinf$
or equivalently $\alpha$).
The fine-structure constant does not "emerge" from new physics; it
is algebraically equivalent to $\Rinf$ and its appearance is a
consistency check.
The non-trivial content is in Postulate~\ref{post:compton}, which
motivates $\rc = r_e/2$ geometrically and provides a self-consistent
classical-hydrodynamic interpretation for all derived scales.

The derivation order matters for non-circularity: the vortex line
tension $\Tc$ must be defined via the Compton-energy
condition~\eqref{eq:Tc_def} before the Laplace-pressure
balance~\eqref{eq:pressure_balance} is applied to obtain
$\rhocore$.
Had $\Tc$ been defined via the Saffman formula~\eqref{eq:Tc_saffman}
from $\rhocore$, the balance would have been vacuous.
The Saffman formula serves only as a downstream consistency check.

The background fluid density $\rhof$ remains open.
Its determination requires an additional independent constraint on
the EM-vortex coupling; the present framework does not supply one.

This algebraic structure arose in the context of a toy model for
topological vortex excitations of the vacuum medium, but the
derivations of Secs.~\ref{sec:derivation}--\ref{sec:derived} are
model-independent: they apply to any vortex-core description
satisfying Postulate~\ref{post:compton}.

% =============================================================
\section{Conclusion}
\label{sec:conclusion}
% =============================================================

We have derived, from Postulate~\ref{post:compton} and the Rydberg
spectroscopic anchor, the complete set of mechanical scales:
\begin{align*}
  \vnorm  &= \sqrt{\pi\hbar c\Rinf/m_e},
  &\rc    &= \vnorm/\omC = r_e/2, \\
  \Gz     &= 2\pi\rc\vnorm,
  &\Tc    &= m_ec^2/(2\rc),\\
  \rhocore &= m_ec^2/(2\pi\vnorm^2\rc^3).
\end{align*}
Each scale has a classical-hydrodynamic interpretation.
The role of $\Rinf$ as input is algebraically equivalent to using
$\alpha$; the derivation makes this equivalence explicit and locates
the physical content in the Compton-core geometric postulate.
Two parameter-free falsifiable predictions are stated.

% =============================================================
\begin{acknowledgments}
The author thanks the referees for identifying the circular
argument in the $\alpha$-claim and the tautological Laplace-balance
in a previous version, leading to the corrected formulation
presented here.
\end{acknowledgments}

% =============================================================
\begin{thebibliography}{99}

\bibitem{Helmholtz1867}
H.~von Helmholtz,
On the integrals of the hydrodynamical equations which express
vortex-motion,
\textit{Philosophical Magazine} \textbf{33}, 485 (1867).

\bibitem{Kelvin1867}
Lord Kelvin (W.~Thomson),
On vortex atoms,
\textit{Proceedings of the Royal Society of Edinburgh} \textbf{6},
94 (1867).

\bibitem{MohrTaylor2018}
P.~J. Mohr, D.~B. Newell, B.~N. Taylor, and E.~Tiesinga,
CODATA recommended values of the fundamental physical constants: 2018,
\textit{Reviews of Modern Physics} \textbf{93}, 025010 (2021).
\doi{10.1103/RevModPhys.93.025010}

\bibitem{Donnelly1991}
R.~J. Donnelly,
\textit{Quantized Vortices in Helium II},
Cambridge University Press, 1991.

\bibitem{Pethick2008}
C.~J. Pethick and H.~Smith,
\textit{Bose--Einstein Condensation in Dilute Gases}, 2nd ed.,
Cambridge University Press, 2008.
\doi{10.1017/CBO9780511802850}

\bibitem{FaddeevNiemi1997}
L.~D. Faddeev and A.~J. Niemi,
Stable knot-like structures in classical field theory,
\textit{Nature} \textbf{387}, 58 (1997).
\doi{10.1038/387058a0}

\bibitem{Barcelo2005}
C.~Barcel\'o, S.~Liberati, and M.~Visser,
Analogue gravity,
\textit{Living Reviews in Relativity} \textbf{8}, 12 (2005).
\doi{10.12942/lrr-2005-12}

\bibitem{Volovik2003}
G.~E. Volovik,
\textit{The Universe in a Helium Droplet},
Oxford University Press, 2003.

\bibitem{Saffman1992}
P.~G. Saffman,
\textit{Vortex Dynamics},
Cambridge University Press, 1992.
\doi{10.1017/CBO9780511624063}

\bibitem{Onsager1949}
L.~Onsager,
Statistical hydrodynamics,
\textit{Nuovo Cimento Supplement} \textbf{6}, 279 (1949).
\doi{10.1007/BF02780991}

\bibitem{Feynman1955}
R.~P. Feynman,
Application of quantum mechanics to liquid helium,
in \textit{Progress in Low Temperature Physics}, vol.~1,
ed.\ C.~J. Gorter, p.~17, North-Holland, 1955.

\bibitem{Gibbons2002}
G.~W. Gibbons,
The maximum tension principle in general relativity,
\textit{Foundations of Physics} \textbf{32}, 1891 (2002).
\doi{10.1023/A:1022370717626}

\bibitem{Schiller2006}
C.~Schiller,
Simple derivation of minimum length, minimum dipole moment, and
lack of space-time continuity,
\textit{International Journal of Theoretical Physics} \textbf{45},
221 (2006).
\doi{10.1007/s10773-005-4835-2}

\bibitem{Karshenboim2005}
S.~G. Karshenboim,
Precision physics of simple atoms: QED tests, nuclear structure, and
fundamental constants,
\textit{Physics Reports} \textbf{422}, 1 (2005).
\doi{10.1016/j.physrep.2005.08.008}

\bibitem{Bethe1947}
H.~A. Bethe,
The electromagnetic shift of energy levels,
\textit{Physical Review} \textbf{72}, 339 (1947).
\doi{10.1103/PhysRev.72.339}

\bibitem{Batchelor1967}
G.~K. Batchelor,
\textit{An Introduction to Fluid Dynamics},
Cambridge University Press, 1967.
\doi{10.1017/CBO9780511800955}

\bibitem{PDG2022}
R.~L. Workman et al.\ (Particle Data Group),
Review of particle physics,
\textit{Progress of Theoretical and Experimental Physics}
\textbf{2022}, 083C01 (2022).
\doi{10.1093/ptep/ptac097}

\end{thebibliography}

\end{document}
